\documentclass[a4paper, 12pt]{report}
\usepackage[utf8]{inputenc}
\usepackage[T1]{fontenc}
\usepackage[french]{babel}
\usepackage{appendix}
\usepackage{graphicx}
\usepackage[a4paper, margin = 2.5 cm]{geometry}
\usepackage{amsmath}

\begin{document}

\begin{center}
  \Huge {\bf UBL} \\
    \vspace{0.5 cm}

  \large Université Libre de Bruxelles\\
  \large Faculté des Sciences Appliquées \\
  \vspace{0.5 cm}
  Année académique 2002-2003\\
    \vspace{7.5 cm}
  \Large Technique de résolution numérique de l'équation de Schrödinger dépendant du temps\\
  \end{center}
  
  \vspace{10 cm}

  \begin{minipage}{0,5 \textwidth}
    \begin{flushleft}
    Directeur de mémoire : \\
    Daniel Baye
  \end{flushleft}
  \end{minipage} \hfill \begin{minipage}{0.5 \textwidth}
    \begin{flushright}
      Travail de fin d'étude \\
      présenté par Gérard Goldstein\\
      en vue de l'obtention du grade\\
      d'Ingénieur Civil Physicien
    \end{flushright}
  \end{minipage}

  \newpage
  
  \vspace*{\fill}
  \begin{flushright}
    \it Je tiens à remercier Daniel Baye et Pierre Capel de m'avoir si bien encadré tout au long de la réalisation de ce travail. Leur grande disponibilité et leur esprit de rigueur ont sans conteste constitué pour moi un formidable atout. Un grand merci également à mes parents et à mon frère pour leur soutien, ainsi qu'à Emma et Julien pour leur solidarité.
  \end{flushright}
  \vspace*{\fill}

  \newpage

  \tableofcontents

  \newpage
  \setcounter{page}{5}
  \chapter*{Introduction}
  \addcontentsline{toc}{chapter}{Introduction}
  L'équation de Schrödingerest l'équation fondamentale de la mécanique quantique. Il s'agit d'une équation aux dérivées partielles qui décrit l'évolution au cours du temps de la fonction d'onde $\Psi$ d'un système physique. Elle prend la forme suivante
  $$i\hbar\frac{\partial \Psi}{\partial t} = H\psi$$
  où H est l'opéarteur hamiltonien(associé à l'énergie totale) du système considéré. Celui-ci
  est la somme des opérateurs d'énergie cinétique et potentielle : 
  $$H(t) = T + V(t)$$
  L’opérateur hamiltonien dépend donc du temps si les potentiels qui entrent en jeu
dépendent eux-mêmes explicitement du temps. Lorsque l’opérateur $H$ ne dépend
pas du temps, on est ramené par séparation des variables spatiales et temporelle
à une équation aux valeurs propres, appelée équation de Schrödinger stationnaire ~\cite{[1]}.Par contre si l’hamiltonien est fonction du temps, on est obligé de résoudre
l'équation de Schrödinger dépendant du temps. C’est le cas par exemple lorsqu’on
traite certains problèmes de manière semi-classique ~\cite{[2], [3]} ou quand on étudie l’effet
sur un atome ou une molécule d’un champ électrique extérieur variable ~\cite{[4]}. Il est
également parfois intéressant de résoudre l’équation de Schrödinger dépendant du
temps même si les potentiels ne dépendent pas du temps. On peut citer entre autres les
problèmes de calcul de coefficients de transmission au travers de barrières de potentiel ~\cite{[5]}. \\

On peut montrer, d’après les propriétés de l’équation de Schrödinger, qu’il existe un
opérateur, appelé opérateur d’évolution qui, appliqué à l’état d’un système à un instant
donné, fournit l’état du système à un instant ultérieur ~\cite{[1]}. La résolution de l’équation de
Schrödinger est alors ramenée à la recherche de cet opérateur qui est lui-même solution
d’une équation, appelée équation d’évolution.

\chapter{L’opérateur d’évolution}
\section{L’équation de Schrödinger}

La mécanique quantique postule qu’à un instant $t_0$ fixé, l’état d’un système physique
est défini par la donné d’un ket (vecteur d’état) $|\psi(t_0)\rangle $ appartenant à l’espace des états $\displaystyle \varepsilon$.\\
En outre l’évolution dans le temps du vecteur d’état $|\psi(t)\rangle $ est régie par l’équation
de Schrödinger

 \begin{equation}
  i\hbar\frac{d}{dt}|\psi(t)\rangle = H(t)|\psi(t)\rangle
  \label{eq:schrodinger}
 \end{equation}
où $H(t)$ est l’observable associée à l’énergie totale du système (l’opérateur hamiltonien
du système). Celle-ci peut dépendre explicitement du temps. \\

\noindent
L’équation de Schrödinger est une équation différentielle du premier ordre par
rapport au temps. Ce qui signifie que la donnée d’un état initial $|\psi(t_0)\rangle$ suffit à déterminer $|\psi(t)\rangle$ à tout instant ultérieur $t$. Ceci n’est valable que si l’évolution n’est
pas interrompue par une mesure d’une grandeur physique du système. \\
Cette équation est également linéaire et homogène. Ses solutions sont donc linéairement
superposables. Si $|\psi_1(t)\rangle$ et $|\psi_2(t)\rangle$ sont deux solutions de l’équation \eqref{eq:schrodinger} et si l’état
initial du système est défini par $|\psi(t_0)\rangle = \lambda_1|\psi_1(t_0)\rangle + \lambda_2|\psi_2(t_0)\rangle $, alors l’état du
système au temps $t$ est donné par $|\psi(t)\rangle = \lambda_1|\psi_1(t)\rangle + \lambda_2|\psi_2(t)\rangle $ . Il existe donc une
correspondance linéaire entre $|\psi(t_0)\rangle$ et $|\psi{t}\rangle $.

\section{Propriétés de l’opérateur d’évolution}

Du fait de la correspondance linéaire entre $|\psi(t_0)\rangle$ et $|\psi{t}\rangle $ , il existe un opérateur
linéaire $U(t,t_0)$ tel que 

\begin{equation}
  \psi(t)\rangle = U(t,t_0)|\psi(t_0)\rangle
  \label{eq:operateur}
\end{equation}

En substituant dans l’équation de Schrödinger \eqref{eq:schrodinger} et \eqref{eq:operateur}, on obtient

$$i\hbar\frac{\partial}{\partial t}U(t,t_0)|\psi(t_0)\rangle = H(t)U(t,t_0)|\psi(t_0)\rangle $$

$\psi(t)\rangle$ pouvant être quelconque, on en déduit l’équation d’évolution

\begin{equation}
  i\hbar\frac{\partial}{\partial t}U(t,t_0) = H(t)U(t,t_0)
  \label{eq:évolution}
\end{equation}

En évaluant l’expression \eqref{eq:operateur} au temps $t = t_0 $ , il advient
$$|\psi(t)\rangle = U(t,t_0)|\psi(t_0)\rangle $$ 
et toujours en considérant que $|\psi(t_0)\rangle$ est quelconque , on obtient le résultat suivant 
label
\begin{equation}
  U(t_0,t_0) = 1
  \label{eq: identité}
\end{equation}
En remplaçant le couple $(t,t_0 )$ dans la définition \eqref{eq:operateur} successivement par les couples $(t,t')$ et $(t' , t'')$ on a 
$$|\psi(t)\rangle = U(t,t')|\psi(t') \quad \mbox{ et } \quad |\psi(t')\rangle = U(t',t'')|\psi(t'')\rangle $$
donc
$$|\psi(t)\rangle = U(t,t')U(t',t'')|\psi(t'')\rangle $$
Mais d'autre part $|\psi(t)\rangle = U(t,t'')|\psi(t'')\rangle $. En comparant ces deux dernières expressions,
on trouve ($|\psi(t'')\rangle $ étant quelconque)
$$ U(t,t'') = U(t,t')U(t',t'') $$
En généralisant :

\begin{equation}
  U(t_n,t_1) = U(t_n, t_{n-1}) \cdots U(t_3, t_2)U(t_2, t_1)
  \label{vecteurdetat}
\end{equation}
On peut également montrer en utilisant la propriété de conservation de la norme 
du vecteur d'état au cours du temps que l'op"rateur d'évolution $U(t,t_0) $ est unitaire. \\
Dans d'autres propriétés de l'équation de Schrödinger et de l'opérateur d'évolution peuvent être trouvées dans ~\cite{[1]} et ~\cite{[11]} 
\end{document}